\chapter{What the virtual laboratory establishes}
\label{ch:vl-positioning}

The virtual laboratory tests whether the inference procedure works when records
are generated through a physical model and an imperfect measurement process.
It separates the hidden process, apparatus and statistical analysis. This makes
it possible to locate a failure: a detector can erase a signal, a clock can
change the model class, a reset can leave dependence between trials, or a
confidence calculation can overstate the evidence. Each calls for a different
correction.

The results have two levels. The proofs specify conditions under which a
certificate is valid. The simulations measure the operating characteristics
of particular implementations under specified data laws. Agreement between an
independent simulator and an analytical prediction verifies the numerical
model. Establishing that an apparatus obeys that model requires physical
measurements. This distinction follows established verification and validation
practice \cite{vlOberkampf2002}.

\section{The experimental contract includes the apparatus}

A temporal trace is generated jointly by dynamics, preparation, detector and
clock. A fixed imperfect detector can often be included in an unknown response
function. A detector that changes with the selected waiting time generally
cannot. Fixed preparation errors can similarly be absorbed into an unknown
initial state, whereas an unrecorded preparation variable correlated with the
controls can enlarge the hidden state space.

This distinction makes a useful robustness question precise: does an
imperfection leave the null class unchanged, merely weaken detection power,
or change the class the theorem describes? The laboratory includes examples
of all three. A counterfeit signal produced by setting-dependent readout is
a demonstration of a missing apparatus assumption. It is not evidence against
a theorem that explicitly requires a fixed detector.

Practical dimension witnesses and processor characterization already have a
substantial experimental literature. Hendrych and collaborators demonstrated
preparation--measurement dimension witnesses \cite{vlHendrych2012}. Gate set
tomography jointly characterizes states, operations and measurement
\cite{vlNielsen2021}. Proctor and collaborators detected and tracked drift in
processor experiments \cite{vlProctor2020}. These results support checking
apparatus stability and retaining acquisition timestamps. They do not establish
those conditions for the illustrative apparatus used here.

The rates, decoherence ratios and readout errors in this laboratory are
\emph{illustrative parameters}. They are not calibrated performance figures
for a superconducting, trapped-ion, photonic or spin platform. A physical
implementation would need joint evidence for its coherence, drive stability,
reset, detector and clock under the intended protocol. A favorable ratio of
frequency to damping alone is insufficient.

\section{Finite-data dissipation inference has close precedents}

Hidden dissipation can already be inferred from apparently time-symmetric
observations. Skinner and Dunkel established such bounds through optimization
over compatible hidden models and applied their framework to biological data
\cite{vlSkinner2021}. Waiting-time methods supply related information
\cite{vlSkinnerWaiting2021}. Finite-sample uncertainty in kinetic inference has
also been treated directly \cite{vlBebon2023}.

The effects of imperfect resolution and finite statistics are particularly
close to the study by Fritz, Ertel and Seifert \cite{vlFritz2025}. Their
comparisons show that finite-data estimates can exceed the underlying entropy
production and that finer resolution need not improve performance at a fixed
observation duration. The practical issue addressed here is therefore a
certified lower confidence bound under an explicit observation contract.

For fresh binary endpoint records, the single-lag floor, monotonicity
constraint, full reversible moment set and predeclared linear contrast use
the same accessible data. Comparisons among them can match record count,
confidence and exposure. A waiting-time estimator requires observed events
that this contract does not supply. Full tomography likewise asks a different
inference question with a richer menu of operations. Neither is a fair
performance baseline if its required observations are withheld.

Combining observables to tighten thermodynamic bounds is established
\cite{vlDechant2021}. The specific improvement studied here compiles a
polynomially certified equilibrium constraint into a single contrast with a
finite-count guarantee. Its validity depends on selecting the contrast before
the held-out records are examined. Optimizing it on those same records requires
an additional selection correction.

\section{Timing errors have different mathematical effects}

Positive exponential mixtures are a standard consequence of reversible
spectral theory \cite{vlAldousFill2002}. Random-time semigroups and their
functional calculus belong to classical subordination theory
\cite{vlSchilling2012}. The timing results here apply those structures to the
measurement contract.

For independent, identically distributed nonnegative time increments $X_j$,
\[
 \mathbb E\!\left[e^{Q(X_1+\cdots+X_k)}\right]
 =\left(\mathbb E[e^{QX_1}]\right)^k.
\]
The same finite hidden state space therefore supports a common stochastic
channel. Both the repeated-channel dimension certificate and the reversible
spectral restrictions survive this type of timing uncertainty.

A common random scale, $T_k=kV$, has a different consequence. For each
reversible spectral mode, $\mathbb E[e^{-\lambda kV}]$ is the $k$th moment of
$x=e^{-\lambda V}$ on $[0,1]$. It preserves the equilibrium moment set. It need
not preserve the same finite-state channel representation: mixing two clock
speeds can turn one decaying mode into two. The clock contributes unobserved
memory if its realization is not retained.

Arbitrary setting-dependent timing laws do not inherit either guarantee
merely because their mean is correct or their errors are bounded. The relevant
condition is the entire family of Laplace profiles used in the theorem. This
provides an apparatus design rule: specify the joint timing mechanism before
assigning it a statistical error bar.

The denominator in a dissipation-rate certificate must also describe physical
exposure. Expected actual waiting time, nominal time and a maximum duration
are different resource statements. An unbounded clock with finite mean has no
pathwise duration cap. Ordinary integration over a detector exposure window
also requires its own observation model; it cannot automatically be replaced
by a randomly selected endpoint.

\section{Dependence is handled through a declared reset guarantee}

Independent trials make binomial intervals appropriate
\cite{vlClopper1934}. Chronological records can instead remain dependent after
an imperfect reset. The dependent-record extension uses a uniform bound on the
conditional bias of each bounded observation relative to its target
probability. Conditional moment-generating-function bounds then control the
fluctuations around those conditional means, and the reset bias enters as a
separate deterministic allowance.

This is an application of established concentration machinery
\cite{caHoeffding1963,vlHoward2021}; mixing-based Markov concentration also has
an extensive theory \cite{vlPaulin2015}. The physical content lies in justifying
the conditional-bias bound. A known worst-case mixing or minorization bound can
supply it after a specified blind waiting interval. A fitted correlation time
from a short record does not establish that uniform condition.
Chapter~\ref{ch:ci-dependence} supplies a separate common-block argument
under a calibrated hidden-process contraction bound.

The supplied argument specifies predictable settings and fixed per-setting
quotas. It does not automatically cover unrestricted data-dependent stopping.
Chapter~\ref{ci:sequential} implements a separate anytime construction
using established confidence-sequence ideas~\cite{vlHoward2021}, conditional
mean bounds and predictable acquisition. It does not license optional
stopping for this fixed-quota result.

\section{Matched design and reproducible evaluation}

The simulation plan separates design data from evaluation data, fixes the
analysis before held-out trials, and retains unsuccessful regimes. This follows
standard simulation methodology \cite{vlMorris2019}. The report should identify
the generated model, accessible records, target quantity, competing methods,
performance measures and random seeds. Independent implementations should
cross-check physical propagation and statistical evaluation where possible.

A scalar contrast can reduce the loss from demanding simultaneous confidence
for every coordinate. Allocating samples in proportion to coefficient size,
and adjusting by the square root of per-record cost, follows classical optimal
stratified allocation \cite{vlCochran1977}. Its use here provides a direct
connection between the chosen physical witness and the measurement budget;
the allocation principle itself is established.

Finite simulation repetitions have uncertainty of their own. If no failures
occur in $R$ independent repetitions, the one-sided $95\%$ exact binomial upper
limit on the failure probability is $1-0.05^{1/R}$. A reported zero count must
therefore retain that uncertainty. Power comparisons on common records should
use paired outcomes, while comparisons with different acquisition menus must
state the additional access and exposure.

The contribution of this laboratory is an auditable bridge from model theorem
to executable protocol: it identifies which imperfections preserve a
certificate, quantifies the resulting confidence and power, and gives explicit
counterexamples when an assumption is removed. Physical validation, external
review and historical priority remain separate stages of evidence.
